\subsection{Retention, estimator and the common cohort}
The whole-archive benchmark comprises \WARepetitions{} station/repetition/setting evaluations at \WAStations{} long-record stations in 18 connected groups. Figure~\ref{fig:whole-archive} and Table~\ref{tab:whole-archive-paired} use common station/repetition rows. For the forward controlled-gap retention contrast, both retention sets are available in \CCCohortCommon{} of \CCCohortAll{} settings; the all-year D2 RMSE on those \CCCohortCommon{} rows is \CCAllCommonRMSE{} $^{\circ}$C per decade, compared with \CCWarmRMSE{} under the warm-season rule. The 820-row all-year value is not used for this paired contrast. Retaining years improves the finite-record slope comparison in the evaluated settings. The executed fixed-budget D2 and classical comparisons vary with time setting and region; small differences do not establish a ranking for more fully sampled source-calibrated annual distributions. Results are conditional on donor availability and mask design. Twenty masks at one station are not twenty independent rivers or climates.

The four-cell factorial comparison gives an arithmetic retention contrast. It is not an orthogonal year-set/reconstruction error decomposition or an identified year-count effect. For offline controlled gaps, the symmetric retention contrast is \WARetentionShare\% of the D2-versus-rule/plug-in MSE gain. This is a result for this experiment, not a universal percentage. Estimator effects and interactions are retained in the full paired table.

The peak-visibility stratification separates cases where the complete reference peak is already visible at the observed lower bound from cases where it is not (Table~\ref{tab:peak-visibility-main}). In the forward empirical-calendar common cases, \CCCalendarVisibleN{} of 3,674 reference peaks are already at the bound and \CCCalendarHiddenN{} are above it; hidden-peak coverage is \CCCalendarHiddenDTwoCov\% for D2 and \CCCalendarHiddenStochCov\% for the classical method (Supporting Information). The controlled-gap stratum is harder. In the forward controlled gaps, \CCPeakVisibleN{} common cases have the reference peak at the bound and \CCPeakHiddenN{} have a higher reference peak. D2 nominal 90\% coverage is \CCPeakVisibleDTwoCov\% and \CCPeakVisibleStochCov\% in the visible stratum, but \CCPeakHiddenDTwoCov\% and \CCPeakHiddenStochCov\% in the hidden stratum. Hidden-stratum MAE is \CCPeakHiddenDTwoMAE{} and \CCPeakHiddenStochMAE{} $^{\circ}$C. Longer gaps, partial training-peak overlap and region change these values: D2 coverage falls from \CCGapThirtyCov\% at 30-day gaps to \CCGapSixtyCov\% at 31--60 days and \CCGapNinetyCov\% at 61--90 days, and from \CCOverlapNoneCov\% with no training-peak overlap to \CCOverlapPartialCov\% with partial overlap. Bavaria has \CCRegionBavariaCov\% coverage versus \CCRegionSwitzerlandCov\% for the Swiss FOEN stations. Removing Bavaria changes the D2-minus-classical MAE difference from \CCOverallMAEDiff{} to \CCLeaveBavariaMAEDiff{} $^{\circ}$C. These strata are descriptive; the hidden-peak stratum uses the withheld reference and cannot become a deployment filter.

Annual marginals and trend uncertainty give different answers (Figure~\ref{fig:joint-uncertainty}). The left panel is a paired comparison on the two-method intersection, weighted by incomplete years. For forward controlled gaps that intersection has \CCPairFwdN{} archives and \CCPairFwdYears{} missing-year cases; annual coverage is \CCPairFwdDTwoAnnual\% for D2 and \CCPairFwdClassicAnnual\% for the classical method. D2 is also available on \CCOwnFwdDTwoN{} of \CCOwnFwdAll{} archives and the classical method on \CCOwnFwdClassicN{}; those availability counts are not the paired denominator. The right panel keeps each method on its own available archives, so the printed sample sizes differ and the panel is not a paired comparison. On D2's own forward set, nominal 95\% independent-rank trend coverage is \WAForwardCoverage\% and residual blocks give \WABlockForwardCoverage\%. On the intersection, independent-rank trend coverage is \CCPairFwdDTwoTrend\% for D2 and \CCPairFwdClassicTrend\% for the classical method. The dashed lines are the nominal 90\% annual target on the left and the nominal 95\% trend target on the right. Within one estimator, the three rank constructions keep that estimator's annual marginals fixed; the two estimators do not share one annual distribution. The joint-rank numerical check increases the samples through 256, 1,024 and \CCRankDTwoDraws{} (Table~\ref{tab:rank-convergence}). At \CCRankDTwoDraws{} samples, nominal 95\% trend coverage is \CCRankStochCovLow\%--\CCRankDTwoCovHigh\% across methods, constructions and seeds, so the undercoverage remains. The largest change in the three-seed mean coverage is \CCRankChangeMax{} percentage points, and the largest single-seed change is \CCRankSeedChangeMax{} percentage points. \CCRankStable{} of \CCRankTotal{} method--construction pairs meet a 0.3-point rule on the three-seed mean. That mean can cancel opposing seed changes, so it is not a resolution guarantee and the constructions are not ranked. The experiment does not uniquely attribute the residual undercoverage to annual miscalibration, tail extrapolation or dependence.

Task calibration is available for \WATaskForwardN{} of 820 forward controlled-gap D2 cases (\WATaskForwardPercent\%), representing two stations and two groups. Both task calibration and the independent-draw comparator cover all of these same cases, but the task-calibrated average width is \CCTaskWidthCal{} $^{\circ}$C per decade versus \CCTaskWidthBase{} for the independent interval. This selected-subset result does not demonstrate general improvement. Withholding and regional support remain explicit.

The three-output release comparison also does not establish a general safety gain. Matching the numerical output count on future long gaps gives \CCReleaseFutureN{} values per rule; the three-output rule produces \CCReleaseThreeError{} errors (\CCReleaseThreeRate\%) and simple score ranking produces \CCReleaseScoreError{} (\CCReleaseScoreRate\%). Matching within the \CCArchiveMatchedN{} forward controlled-gap archives gives trend RMSE \CCArchiveMatchedThreeRMSE{} versus \CCArchiveMatchedScoreRMSE{} $^{\circ}$C per decade, and the connected-group interval for the paired squared-error difference spans zero. The published-values-plus-bounds working likelihood treats released values as exact and covers 99.8\% with a 3.26 $^{\circ}$C per decade width; it is not a calibrated missingness-only interval. These are boundary results, not a deployment recommendation.
